\chapter{R, RStudio, Git, and Reproducible Project Structure}

\WeekHeader{2}{R, RStudio, Git, and Reproducible Project Structure}{Build a project that another researcher can open, run, and audit.}{Working R project with folders, install check, research log, and first commit.}

\section{Why This Matters}
Reproducibility is a research skill, not an afterthought. A student who can organize files, name outputs clearly, and rerun an analysis has already solved many of the problems that slow down publication.

\section{Concepts}
\begin{itemize}
  \item \textbf{Project root}: the top-level folder that contains scripts, data, results, and documentation.
  \item \textbf{Relative path}: a file path that works from the project root on any machine.
  \item \textbf{Raw data}: input data that should not be manually edited.
  \item \textbf{Derived data}: tables created by scripts.
  \item \textbf{Version control}: a record of changes to files and analysis decisions.
  \item \textbf{Research log}: a dated record of decisions, errors, fixes, and interpretation changes.
\end{itemize}

\section{Recommended Project Structure}
\begin{verbatim}
student-project/
  README.md
  analysis/
  data_raw/
  data_processed/
  figures/
  results/
  scripts/
  reports/
  logs/
\end{verbatim}

\section{Guided Lab}
Run the install and environment check in \path{training/scripts/00_install_check.R}. Students should save the output in their research log.

\begin{rcode}
source("training/scripts/00_install_check.R")
check_training_install()
\end{rcode}

Then create a project folder and initialize the file structure:

\begin{rcode}
source("training/scripts/01_project_setup.R")
create_training_project("student-project")
\end{rcode}

\section{Git Practice}
Students should make small, readable commits. A good commit message says what changed and why it matters:

\begin{verbatim}
Add chemical intake table for plant volatile project
Document first uafR install check
Add Week 1 project brief
\end{verbatim}

\section{Independent Extension}
Students customize their project README with:
\begin{itemize}
  \item Project title.
  \item Research question.
  \item Data source.
  \item Main scripts.
  \item Contact and dsDNA Core support.
  \item Current status.
\end{itemize}

\section{Deliverables}
\begin{tabularx}{\textwidth}{p{0.25\textwidth}YY}
\DeliverableTableHeader
Project folder & Uses the required folder structure and relative paths. & Local folder or repository\\
Install check & Shows R version, uafR version, key dependencies, and package data availability. & Research log\\
First commit & Commit includes Week 1 and Week 2 materials without private tokens or sensitive data. & Git history\\
\end{tabularx}

\section{Quality Checklist}
\begin{checklistbox}{Before Week 3}
\begin{itemize}
  \item No script uses an absolute path to another person's computer.
  \item No token, password, or private key appears in any file.
  \item The project can be opened from the project root and scripts can find their inputs.
  \item The student can explain the difference between raw and processed data.
\end{itemize}
\end{checklistbox}
